\DocumentMetadata{
  pdfversion=2.0,
  pdfstandard=ua-2,
  lang=en-US,
  tagging=on,
  tagging-setup={math/setup=mathml-SE,tabsorder=structure}
}
\documentclass[11pt,letterpaper]{article}
\usepackage[margin=0.68in,includefoot]{geometry}
\usepackage{newcomputermodern}
\usepackage{amsmath,amscd,mathtools}
\providecommand{\square}{\mdlgwhtsquare}
\usepackage[hidelinks]{hyperref}
\hypersetup{
  pdftitle={Analysis First-Year Exam, May 2024, Part 2},
  pdfauthor={University of Florida Department of Mathematics},
  pdfsubject={Graduate mathematics examination},
  pdfdisplaydoctitle=true,
  bookmarksopen=true
}
\setcounter{secnumdepth}{0}
\setlength{\parindent}{0pt}
\setlength{\parskip}{0.28em}
\setlength{\emergencystretch}{3em}
\setlength{\leftmargini}{2.15em}
\setlength{\leftmarginii}{2.65em}
\setlength{\leftmarginiii}{3em}
\renewcommand{\labelenumi}{\textbf{\theenumi.}}
\pagestyle{empty}
\raggedbottom
\allowdisplaybreaks
\newenvironment{examproblems}[1]{%
  \begin{enumerate}%
  \setcounter{enumi}{#1}%
  \setlength{\topsep}{0.35em}%
  \setlength{\partopsep}{0pt}%
  \setlength{\itemsep}{0.48em}%
  \setlength{\parsep}{0.18em}%
}{\end{enumerate}}
\newenvironment{examsubparts}{%
  \begin{enumerate}%
  \setlength{\topsep}{0.18em}%
  \setlength{\partopsep}{0pt}%
  \setlength{\itemsep}{0.16em}%
  \setlength{\parsep}{0.1em}%
}{\end{enumerate}}
\newenvironment{examinstructions}{%
  \par\begingroup\itshape
}{%
  \par\endgroup\vspace{0.18em}
}
\makeatletter
\renewcommand\section{\@startsection{section}{1}{0pt}%
  {-1.25ex plus -.25ex minus -.1ex}%
  {0.55ex plus .1ex}%
  {\normalfont\large\bfseries}}
\renewcommand{\@maketitle}{%
  \newpage\null\vskip -1.2em%
  {\footnotesize\bfseries
    \href{https://www.ufl.edu/}{University of Florida}, \href{https://math.ufl.edu/}{Department of Mathematics}\par}%
  \vskip 0.18em%
  \tagstructbegin{tag=Title}%
    {\LARGE\bfseries\href{https://gma.math.ufl.edu/exams/analysis-first-year/}{Analysis First-Year Exam}, May 2024, Part 2\par}%
  \tagstructend%
  \vskip 0.35em%
  \tagmcbegin{artifact}%
    \hrule height 1pt%
  \tagmcend%
  \vskip 0.55em%
}
\makeatother
\title{Analysis First-Year Exam, May 2024, Part 2}
\author{}
\date{}
\begin{document}
\maketitle
\thispagestyle{empty}
\begin{examinstructions}
Answer FOUR questions in detail. State carefully any results used without proof.
\end{examinstructions}
\begin{examproblems}{0}
\item
Let \(f:[a,b]\to\mathbb{R}\) be Riemann integrable, let~\(g\) be a real-valued function on \([a,b]\), and let \(D=\{t\in[a,b]:f(t)\ne g(t)\}\).
\begin{examsubparts}
\item[\textnormal{(i)}]
Show by example that if~\(D\) is countably infinite then~\(g\) can fail to be Riemann integrable.
\item[\textnormal{(ii)}]
Prove that if~\(D\) is finite then~\(g\) must be Riemann integrable.
\end{examsubparts}
\item
Prove Dini’s theorem that, if the sequence \((f_n)_{n=0}^{\infty}\) of continuous real-valued functions on a compact space decreases pointwise to zero, then the convergence is uniform. Show by example that uniform convergence can fail if decreases pointwise is replaced by converges pointwise.
\item
State the Weierstrass approximation theorem. Determine exactly which continuous functions \(f:[-1,1]\to\mathbb{R}\) satisfy the requirement that for each strictly positive integer~\(n\), 
\[
\int_{-1}^{1}f(t)t^n\,dt=0.
\]
\item
Let \((\Omega,\mathcal{A},\mu)\) be a measure space and let \((A_n)_{n=1}^{\infty}\) be a sequence in~\(\mathcal{A}\). For each positive integer~\(N\), let \(B_N\) be the set containing exactly those \(\omega\in\Omega\) that lie in \(A_n\) for at least~\(N\) values of~\(n\). By considering the series \(\sum_{n=1}^{\infty}\mathbf{1}_{A_n}\) of indicator functions, prove that 
\[
N\mu(B_N)\leq\sum_{n=1}^{\infty}\mu(A_n).
\]
\item
Let the non-negative function \(f:\mathbb{R}\to\mathbb{R}\) be integrable with respect to Lebesgue measure~\(\lambda\) and for each real~\(t\) define 
\[
F(t)=\int_0^t f\,d\lambda.
\]
 Prove that the function~\(F\) is continuous. Must~\(F\) be uniformly continuous? Explain.
\end{examproblems}
\end{document}
